\chapter{Simulation Methods and Architecture}

\section{Spatial Discretization: Voxel Mechanics}
To manage the computational complexity of the continuous physical world, Layer 1 employs a discrete, voxel-based spatial representation. The environment is partitioned into $32^3$ voxel chunks, an architecture heavily inspired by cellular automata and sparse voxel octrees (SVOs) \cite{laine2010efficient}.

\subsection{Voxel State and Material Properties}
Each voxel encapsulates the localized physical state, rather than relying on continuous meshes. A voxel state tensor includes:
\begin{itemize}
    \item \textbf{Material ID:} Defines the atomic or molecular composition (e.g., \texttt{HARDWARE\_FAB\_NODE}, \texttt{POLYMER\_SPOOL}).
    \item \textbf{Intensive Properties:} Temperature, moisture (relative humidity), and chemical potential.
    \item \textbf{Metadata Bitmask:} Encodes actuator properties, sensor presence, and network states.
\end{itemize}

\subsection{Local Thermal Propagation}
Thermal exchange between voxels is simulated using a discrete approximation of Fourier's law of thermal conduction. For a voxel $i$ adjacent to voxel $j$, the heat transfer rate $\dot{q}$ is evaluated across the discrete boundary:
\begin{equation}
\dot{q}_{i \to j} = \frac{k_{ij} A}{d} (T_i - T_j)
\end{equation}
where $k_{ij}$ is the effective thermal conductivity at the interface, $A$ is the voxel face area, and $d$ is the voxel width. This cellular automata approach ensures that thermodynamic operations (like heat dissipation from an FDM printer) remain localized to the chunk, allowing for highly parallelized $O(1)$ state updates \cite{wolfram1983statistical}.



\section{Acoustic Wave Propagation and Attenuation}
In modeling communal physical spaces (such as the environments described in Scenario Epsilon: Coworking), the simulation must account for the propagation of longitudinal mechanical waves (sound). An "acoustic budget" is not a mere social construct; it is a physical parameter tracked by the voxel engine.

\subsection{Discrete Acoustic Modeling}
Sound propagation through the $32^3$ chunk grid is modeled using a simplified wave equation. For a scalar acoustic pressure field $p(\vec{r}, t)$, the propagation is governed by:
\begin{equation}
\frac{1}{c^2} \frac{\partial^2 p}{\partial t^2} - \nabla^2 p = 0
\end{equation}
where $c$ is the speed of sound in the local voxel medium. To model soundproofing (e.g., a podcast booth), the voxel metadata includes an acoustic impedance factor $Z = \rho c$ (where $\rho$ is density). When an acoustic wave transitions between voxels of differing impedance $Z_1$ and $Z_2$, the transmission coefficient $T$ dictates the attenuation of the wave:
\begin{equation}
T = \frac{4 Z_1 Z_2}{(Z_1 + Z_2)^2}
\end{equation}
This allows the orchestrator to physically verify if a user in a "DeepWork" zone is subjected to acoustic energy exceeding their threshold, dynamically adjusting reputational trust scores if a user violates the physical acoustic budget of the space \cite{kinsler2000fundamentals}.



\subsection{Algorithmic Propagation and Boundary Conditions}
To compute acoustic transmission across large voxel grids efficiently (e.g., Scenario Lambda: Cultural Events and Somatic Play), the engine employs a modified breadth-first search (BFS) algorithm. The sound intensity $I$ at distance $r$ strictly follows the inverse square law in uniform media, but is further attenuated by the cumulative absorption coefficient $\alpha$ of intersected voxels:
\begin{equation}
I(r) = \frac{I_0}{4\pi r^2} e^{-\sum \alpha_i \Delta x_i}
\end{equation}
By continuously mapping $I(r)$ to the boundary layer of the chunk grid (the \texttt{PROPERTY\_LINE}), the simulation can autonomously actuate visual feedback loops (warning lights) if the intensity approaches municipal legal limits, thereby modeling civil acoustic compliance as a rigid boundary value problem.

\section{Network Physics and Percolation Federation}
When individual nodes sever their connections to centralized legacy servers (Scenario Omega), the simulation must route telemetry and state syncs across a purely decentralized physical mesh (e.g., LoRaWAN, dark fiber). The resilience of this Gossipsub federation is modeled using percolation theory.

\subsection{Topological Phase Transitions}
The mesh network is represented as a graph $G(V, E)$, where $V$ are the physical nodes and $E$ are the active P2P connections. As nodes federate, the probability $p$ of a connection existing between any two adjacent nodes dictates the macro-state of the network. According to percolation theory, there exists a critical threshold $p_c$. When $p > p_c$, a topological phase transition occurs, and a giant connected component spans the entire network, ensuring that physical telemetry and epistemic data can propagate globally with near certainty, despite the absence of centralized infrastructure \cite{stauffer2014introduction}.



\subsection{High-Gradient Thermal Diffusion and Ignition}
In high-temperature industrial scenarios (e.g., Scenario Theta: The Foundry), the discrete approximation of Fourier's law must accommodate extreme thermal gradients. When a voxel (such as an induction kiln) reaches $T > 1000^\circ\text{C}$, the radiative heat transfer $\dot{q}_{rad}$ dominates, scaling with the fourth power of temperature via the Stefan-Boltzmann law:
\begin{equation}
\dot{q}_{rad} = \varepsilon \sigma A (T_{kiln}^4 - T_{adj}^4)
\end{equation}
where $\varepsilon$ is emissivity and $\sigma$ is the Stefan-Boltzmann constant. The simulation engine dynamically calculates this radiative flux against the specific ignition threshold of adjacent voxels. If the kiln is placed adjacent to a voxel with low thermal mass and a low ignition point (e.g., \texttt{WOOD\_WALL}), the thermal bleed rapidly forces the adjacent voxel into a \texttt{FIRE} state, mathematically punishing spatial designs that ignore foundational thermodynamics \cite{incropera2007fundamentals}.



\section{Microgrid Islanding and Power Shedding}
In modeling catastrophic shocks to the legacy system (Scenario Upsilon), the Oasis engine must transition the electrical simulation from an infinite-supply model (grid connected) to a closed-system microgrid.

\subsection{Dynamic Load Balancing}
When a grid failure is detected, the simulation strictly enforces the conservation of energy across the local electrical network. The total power consumed $P_{load}$ must be perfectly matched by the local generation $P_{gen}$ and battery discharge $P_{batt}$:
\begin{equation}
P_{load}(t) = P_{gen}(t) + P_{batt}(t)
\end{equation}
Because $P_{gen}$ (e.g., solar) is highly variable and $P_{batt}$ is finite, the BPMN engine acts as an autonomous shedding controller. It iterates through all active voxels, reading their metadata. Voxels with a \texttt{priority\_tier} below the critical threshold are mathematically severed from the power bus ($P_{load, i} \to 0$). This extends the $dt$ (time differential) of the battery lifespan before $P_{batt} \to 0$, ensuring critical survival states (like medical refrigeration) are maintained through the shock \cite{lasseter2002microgrids}.



\section{Spatial Semaphores and Temporal Voxel Locking}
To model the private allocation of space (Scenario Zeta: Lodging), the Oasis engine applies the computer science principles of mutual exclusion (mutex) directly to physical coordinates.

\subsection{Physical Locking Mechanisms}
When a chunk of space is designated for private use (e.g., a micro-lodge or a reserved mediation room from Scenario Psi), the \texttt{ChunkManager} treats the bounding box as a critical section. A spatial semaphore is applied to the voxels. During the temporal block of the occupancy state, any entity lacking the cryptographic key (the active tenant or the neutral mediator) will experience an absolute pathfinding collision against the boundary voxels, even if the material is physically permeable (like a tent). This ensures that cryptographic privacy is perfectly enforced as a rigid physical law within the simulation boundaries \cite{dijkstra1965solution}.

\begin{thebibliography}{99}
\bibitem{kinsler2000fundamentals}
Kinsler, L. E., Frey, A. R., Coppens, A. B., \& Sanders, J. V. (2000). \textit{Fundamentals of acoustics}. John Wiley \& Sons.

\bibitem{incropera2007fundamentals}
Incropera, F. P., Lavine, A. S., Bergman, T. L., \& DeWitt, D. P. (2007). \textit{Fundamentals of heat and mass transfer}. John Wiley \& Sons.

\bibitem{lasseter2002microgrids}
Lasseter, R. H. (2002). \textit{Microgrids}. In 2002 IEEE Power Engineering Society Winter Meeting. Conference Proceedings (Vol. 1, pp. 305-308). IEEE.

\bibitem{stauffer2014introduction}
Stauffer, D., \& Aharony, A. (2014). \textit{Introduction to percolation theory}. Taylor \& Francis.

\bibitem{dijkstra1965solution}
Dijkstra, E. W. (1965). \textit{Solution of a problem in concurrent programming control}. Communications of the ACM, 8(9), 569.

\bibitem{wolfram1983statistical}
Wolfram, S. (1983). \textit{Statistical mechanics of cellular automata}. Reviews of modern physics, 55(3), 601.

\bibitem{laine2010efficient}
Laine, S., \& Karras, T. (2010). \textit{Efficient sparse voxel octrees}. IEEE Transactions on Visualization and Computer Graphics, 17(8), 1048-1059.

\end{thebibliography}
